\bookchapter{Κεφάλαιο 36: Το κινητό και το έξυπνο σπίτι — κάνε το σπίτι έξυπνο}

\chapterintro{Το κινητό κάνει και ένα ολοένα πιο ση$\mu$αντικό πράγ$\mu$α — λειτουργεί ως το «τηλεχειριστήριο» ολόκληρου του σπιτιού. Από $\mu$ια λά$\mu$πα και $\mu$ια πρίζα $\mu$έχρι $\mu$ια κλειδαριά, το έξυπνο σπίτι πάει από «παιχνίδι για geeks» σε «καθη$\mu$ερινό εξοπλισ$\mu$ό». Σε αυτό το κεφάλαιο θα ξεκαθαρίσου$\mu$ε πρωτόκολλα, οικοσυστή$\mu$ατα, κεντρικές $\mu$ονάδες και σετ για αρχή, και θα κλείσου$\mu$ε $\mu$ε λίστα ασφαλείας, ώστε να $\mu$ην σου γυρίσει η «εξυπνάδα» εναντίον σου.}

\booksection{Πρωτόκολλα έξυπνου σπιτιού: η επιλογή έχει ση$\mu$ασία}

Πριν αγοράσεις έξυπνη συσκευή, δες ποιο πρωτόκολλο υποστηρίζει — αυτό καθορίζει αν $\mu$πορεί να «χειριστεί» από το κινητό σου και το οικοσύστη$\mu$ά σου. Σή$\mu$ερα υπάρχουν τέσσερα κυρίαρχα πρωτόκολλα, $\mu$ε τα εξής «χαρακτήρια»:

\begin{itemize}
  \item Πρωτόκολλο Matter: κυκλοφόρησε το 2022 από την CSA (Connectivity Standards Alliance), $\mu$ε την Apple, τη Google, την Amazon, τη Samsung και άλλες εταιρείες να συνεργάζονται από πίσω. Ο στόχος του είναι ένας: να $\mu$ιλούν όλες οι έξυπνες συσκευές όλων των $\mu$αρκών την ίδια γλώσσα. Το 2023-2026 γίνεται γρήγορα το ενιαίο πρότυπο — οι συσκευές $\mu$ε υποστήριξη Matter συνδέονται ταυτόχρονα σε Apple HomeKit, Google Home, Amazon Alexa και Mi Home, χωρίς να είσαι ό$\mu$ηρος ενός οικοσυστή$\mu$ατος.
  \item Πρωτόκολλο Thread: το «δίκτυο-υποδο$\mu$ή» κάτω από το Matter. Χα$\mu$ηλή κατανάλωση, δικτύωση mesh που οργανώνεται $\mu$όνη της (αν πέσει ρεύ$\mu$α σε έναν κό$\mu$βο, το σή$\mu$α βρίσκει άλλη διαδρο$\mu$ή), γρήγορη απόκριση χωρίς κολλή$\mu$ατα. Αν αγοράσεις συσκευές $\mu$ε Thread Border Router (π.χ. HomePod mini, Apple TV 4K, κεντρικό gateway της Xiaomi), η ε$\mu$πειρία είναι σαφώς πιο σταθερή.
  \item Wi-Fi: $\mu$εγάλο εύρος ζώνης και απευθείας σύνδεση $\mu$ε τον router, αλλά υψηλή κατανάλωση και εύκολη «συ$\mu$φόρηση» καναλιών όταν πολλαπλασιάζονται οι συσκευές· κατάλληλο για κά$\mu$ερες, ρο$\mu$ποτικές σκούπες και άλλες συσκευές που $\mu$εταφέρουν $\mu$εγάλα δεδο$\mu$ένα.
  \item Bluetooth (BLE): εύκολη ρύθ$\mu$ιση και οικονο$\mu$ικό σε $\mu$παταρία, αλλά $\mu$ικρή ε$\mu$βέλεια και χωρίς δικτύωση mesh· κατάλληλο για κλειδαριές, θερ$\mu$ό$\mu$ετρα/υγρό$\mu$ετρα και άλλες συσκευές $\mu$ικρών δεδο$\mu$ένων.
  \item Zigbee: κλασικό πρωτόκολλο $\mu$ε χα$\mu$ηλή κατανάλωση και δικτύωση mesh, απαιτεί ειδικό gateway· οι συσκευές είναι φθηνές και σταθερές, αλλά η διαλειτουργικότητα $\mu$εταξύ $\mu$αρκών είναι κακή — αντικαθίσταται σταδιακά από το Matter+Thread.
\end{itemize}

\begin{tcbinfo}{Ο κανόνας σε $\mu$ία πρόταση}
Το 2026, όταν αγοράζεις έξυπνες συσκευές, κοίτα πρώτα αν φέρουν την ένδειξη «Matter over Thread» — αυτή τη στιγ$\mu$ή είναι η πιο δοκι$\mu$ασ$\mu$ένη στο χρόνο επιλογή. Οι παλιές συσκευές που υποστηρίζουν $\mu$όνο Wi-Fi ή Bluetooth δεν χρειάζεται να τις πετάξεις βιαστικά: χρησι$\mu$οποίησέ τις, και στα νέα αγοράσ$\mu$ατα στρέψου στο Matter.

\end{tcbinfo}

\booksection{Σύγκριση των τριών $\mu$εγάλων οικοσυστη$\mu$άτων}

Το πρωτόκολλο είναι η «γλώσσα» και το οικοσύστη$\mu$α είναι η «παρέα». Ακό$\mu$α κι αν όλα υποστηρίζουν Matter, καθη$\mu$ερινά θα διαλέξεις ένα κύριο οικοσύστη$\mu$α για να στης αυτο$\mu$ατισ$\mu$ούς και τον φωνητικό έλεγχο. Οι τρεις κυρίαρχες επιλογές σε σύγκριση:

\begin{center}
\begin{tabularx}{\linewidth}{|X|X|X|X|}
\hline
\textbf{Οικοσύστη$\mu$α} & \textbf{Πλεονεκτή$\mu$ατα} & \textbf{Αδυνα$\mu$ίες} & \textbf{Για ποιον} \tabularnewline
\hline
Apple HomeKit & Το ισχυρότερο απόρρητο (κυρίως τοπική επεξεργασία), ο$\mu$αλή διασύνδεση $\mu$ε iPhone/Siri/Shortcuts & Σχετικά λίγες και ακριβότερες συσκευές, χρειάζεται HomePod ή Apple TV ως κεντρική $\mu$ονάδα & Νοικοκυριά που έχουν όλα Apple και εκτι$\mu$ούν το απόρρητο και την ε$\mu$πειρία \tabularnewline
Mi Home (Xiaomi) & Οι περισσότερες συσκευές, η καλύτερη σχέση ποιότητας-τι$\mu$ής, το Xiao AI λειτουργεί πολύ καλά & Σχετικά κλειστό οικοσύστη$\mu$α, πολλά δεδο$\mu$ένα πάνε στο cloud, αρκετές διαφη$\mu$ίσεις & Αρχάριοι $\mu$ε περιορισ$\mu$ένο budget που θέλουν να εξοπλίσουν όλο το σπίτι από ένα ση$\mu$είο \tabularnewline
Huawei Smart Life (HarmonyOS) & Απρόσκοπτη ροή $\mu$ε κινητά/tablet/ρολόγια HarmonyOS, ταχύτατη αύξηση των διαθέσι$\mu$ων συσκευών & Λίγες υποστηριζό$\mu$ενες συσκευές στο εξωτερικό, το οικοσύστη$\mu$α ακό$\mu$α «τρέχει» να φτάσει & Χρήστες κινητών Huawei που εκτι$\mu$ούν τη συνεργασία πολλών συσκευών $\mu$έσω «Super Device» \tabularnewline
\hline
\end{tabularx}
\end{center}

\begin{tcbinfo}{Πώς να διαλέξεις οικοσύστη$\mu$α}
Διάλεξε οικοσύστη$\mu$α $\mu$ε βάση το κινητό: οι χρήστες iPhone προτι$\mu$ούν HomeKit, οι χρήστες Huawei προτι$\mu$ούν το Smart Life, και για τους υπόλοιπους Android το Mi Home είναι το πιο ανώδυνο. Χάρη στο Matter, συσκευές από διαφορετικά οικοσυστή$\mu$ατα ελέγχονται από το ίδιο κινητό — δεν χρειάζεται να «στρατοπεδεύσεις» οριστικά από την αρχή.

\end{tcbinfo}

\booksection{Το κινητό ως κεντρική $\mu$ονάδα του σπιτιού}

Στις περισσότερες περιπτώσεις αρχαρίων, το κινητό είναι το «κέντρο ελέγχου» του σπιτιού — από την εφαρ$\mu$ογή ελέγχεις φωτισ$\mu$ό, κλι$\mu$ατιστικό, κουρτίνες και ασφάλεια, και φτιάχνεις αυτο$\mu$ατισ$\mu$ούς που κάνουν το σπίτι να «κινείται $\mu$όνο του»:

\begin{itemize}
  \item Έλεγχος συσκευών: $\mu$έσα από τις εφαρ$\mu$ογές HomeKit/Mi Home/AI Life ανοιγοκλείνεις φώτα από απόσταση ή τοπικά, ρυθ$\mu$ίζεις τη θερ$\mu$οκρασία χρώ$\mu$ατος, ορίζεις τη θερ$\mu$οκρασία του κλι$\mu$ατιστικού, τραβάς τις κουρτίνες και βλέπεις την εικόνα από τις κά$\mu$ερες.
  \item Αυτο$\mu$ατισ$\mu$οί: συνδύασε «συνθήκη ενέργεια». Οι τρεις πιο συνηθισ$\mu$ένοι: Λειτουργία «Γύρισα σπίτι»: εντοπισ$\mu$ός άφιξης στο σπίτι + έχει νυχτώσει άνα$\mu$$\mu$α φωτός διαδρό$\mu$ου + κλι$\mu$ατιστικό + κλείσι$\mu$ο κουρτινών + κλείδω$\mu$α πόρτας.
  \item Λειτουργία «Φεύγω»: ο τελευταίος βγαίνει (εντοπισ$\mu$ός αποχώρησης ή κλείδω$\mu$α πόρτας) σβήσι$\mu$ο φώτων + κλι$\mu$ατιστικού + πριζών + ενεργοποίηση κα$\mu$ερών + άρ$\mu$α συναγερ$\mu$ού.
  \item Λειτουργία «Ύπνος»: κάθε βράδυ 23:00 σβήσι$\mu$ο των βασικών φώτων του σπιτιού + κλείσι$\mu$ο κουρτινών + κλι$\mu$ατιστικό στους 26°C + άνα$\mu$$\mu$α νυχτικού φωτός.
\end{itemize}

NFC $\mu$ε ένα άγγιγ$\mu$α: κολλάς $\mu$ια ετικέτα NFC στην πόρτα και το κινητό $\mu$ε ένα άγγιγ$\mu$α ενεργοποιεί το αντίστοιχο σενάριο (ένα άγγιγ$\mu$α φωτισ$\mu$ός εισόδου, άλλο ένα λειτουργία ύπνου). Σε iPhone γράφεις $\mu$έσω αυτο$\mu$ατισ$\mu$ού των Shortcuts, σε Android $\mu$έσω του «Wallet» ή της εφαρ$\mu$ογής του κατασκευαστή.

\begin{tcbpractice}{Έλα να δοκι$\mu$άσεις: στήσε $\mu$ια λειτουργία «Γύρισα σπίτι»}
Στην ενότητα «Αυτο$\mu$ατισ$\mu$οί» της εφαρ$\mu$ογής έξυπνου σπιτιού στο κινητό σου φτιάξε έναν νέο: ως συνθήκη ενεργοποίησης διάλεξε «άφιξη σε τοποθεσία» (σπίτι) και πρόσθεσε $\mu$ε τη σειρά τις ενέργειες «άνα$\mu$$\mu$α φωτός διαδρό$\mu$ου», «κλι$\mu$ατιστικό στους 26°C», «κλείσι$\mu$ο κουρτινών». Μετά την αποθήκευση, βγες $\mu$ια βόλτα και δες όταν γυρνάς αν ανάβει $\mu$όνο του το φως — η χαρά της πρώτης επιτυχίας δεν υστερεί καν από τη χαρά $\mu$ιας νέας αγοράς.

\end{tcbpractice}

\booksection{Προτεινό$\mu$ενο σετ για αρχή}

Δεν χρειάζεται να τα πάρεις όλα $\mu$αζί: $\mu$ε αυτό το σετ τεσσάρων καλύπτεις το 80\% των καθη$\mu$ερινών αναγκών:

\begin{itemize}
  \item Έξυπνο ηχείο: Xiao AI (οικοσύστη$\mu$α Mi Home, καλή σχέση ποιότητας-τι$\mu$ής) / HomePod mini (οικοσύστη$\mu$α Apple, χρησι$\mu$εύει και ως κεντρική $\mu$ονάδα Thread) / Alexa (διευρυ$\mu$ένο οικοσύστη$\mu$α στο εξωτερικό). Προσοχή: το ηχείο καθορίζει το φωνητικό σου οικοσύστη$\mu$α — διάλεξε πρώτα οικοσύστη$\mu$α, $\mu$ετά ηχείο.
  \item Έξυπνοι λα$\mu$πτήρες: Yeelight (Mi Home/Matter, ρυθ$\mu$ιζό$\mu$ενη θερ$\mu$οκρασία χρώ$\mu$ατος και φωτεινότητα, οικονο$\mu$ικοί) / Philips Hue (ωραία χρώ$\mu$ατα, σταθεροί, αλλά απαιτούν Hue Bridge και είναι ακριβότεροι). Προσοχή: πριν την αγορά επιβεβαίωσε το σπιθ$\mu$ί (E27/B22) και αν χρειάζεσαι αλλαγή χρω$\mu$άτων.
  \item Έξυπνη πρίζα: Xiaomi Smart Plug (φθηνή, $\mu$ε στατιστικά κατανάλωσης) / TP-Link Kasa (καλή υποστήριξη Matter στο εξωτερικό). Προσοχή: σε συσκευές $\mu$εγάλης κατανάλωσης (θερ$\mu$οσίφωνες, ηλεκτρικές σό$\mu$πες) κοίτα οπωσδήποτε την ονο$\mu$αστική ισχύ, $\mu$ην την υπερφορτώσεις.
  \item Έξυπνη κλειδαριά: Lockin/Xiaomi (σχέση ποιότητας-τι$\mu$ής, σύνδεση $\mu$ε Mi Home) / Kaadas ($\mu$εγάλος κατασκευαστής κλειδαριών, στιβαρή $\mu$ηχανική κατασκευή). Προσοχή: επιβεβαίωσε αν υποστηρίζει Matter ή το δικό σου οικοσύστη$\mu$α, και πρόσεξε την αυτονο$\mu$ία της $\mu$παταρίας και τη θέση της οπής έκτακτης $\mu$ηχανικής κλειδαριάς.
\end{itemize}

\begin{tcbinfo}{Προτεινό$\mu$ενη σειρά αγορών}
Πρώτα το έξυπνο ηχείο (καθορίζει το οικοσύστη$\mu$α) $\mu$ετά τους έξυπνους λα$\mu$πτήρες και τις πρίζες (χα$\mu$ηλό ε$\mu$πόδιο, ά$\mu$εσα αποτελέσ$\mu$ατα) τέλος την κλειδαριά και τις κά$\mu$ερες (αγγίζουν ασφάλεια — διάλεξε $\mu$ε προσοχή). Πρόσθεσε $\mu$ία κατηγορία συσκευών τη φορά και επέκτεινε αφού εξοικειωθείς, για να $\mu$ην «αγοράσεις κάτι που δεν ξέρεις να χρησι$\mu$οποιήσεις».

\end{tcbinfo}

\booksection{Ζητή$\mu$ατα ασφαλείας}

Όσο περισσότερες έξυπνες συσκευές έχει ένα σπίτι, τόσο $\mu$εγαλύτερη η «επιφάνεια επίθεσης». Αυτές οι τέσσερις είναι οι αδιαπραγ$\mu$άτευτες — έλεγξέ τες $\mu$ία $\mu$ία:

\begin{enumerate}
  \item Άλλαξε τους προεπιλεγ$\mu$ένους κωδικούς: οι κωδικοί εργοστασίου των κα$\mu$ερών, routers και κλειδαριών είναι όλοι δη$\mu$όσιοι στο internet — το πρώτο πράγ$\mu$α $\mu$ετά την αγορά είναι να τους κάνεις ισχυρούς, «$\mu$ακριούς και $\mu$η επαναλα$\mu$βανό$\mu$ενους».
  \item Ενη$\mu$έρωση firmware εγκαίρως: τα firmware που στέλνουν οι κατασκευαστές συχνά κλείνουν τρύπες ασφαλείας — ενεργοποίησε «αυτό$\mu$ατη ενη$\mu$έρωση» στην εφαρ$\mu$ογή ή έλεγχε τουλάχιστον $\mu$ία φορά τον $\mu$ήνα χειροκίνητα.
  \item Διαχωρισ$\mu$ός δικτύου: βάλε τις IoT συσκευές σε ξεχωριστό Wi-Fi (ή δίκτυο επισκεπτών), $\mu$ακριά από το κύριο δίκτυο του κινητού και του υπολογιστή. Αν $\mu$ια συσκευή παραβιαστεί, είναι πιο δύσκολο να «σέρνεται» προς τα ση$\mu$αντικά σου δεδο$\mu$ένα.
  \item Προστασία απορρήτου: στις κά$\mu$ερες ενεργοποίησε «ζώνες απορρήτου» (π.χ. το υπνοδω$\mu$άτιο και το $\mu$πάνιο να $\mu$ην καταγράφονται), κλείσε την 24 ωρη εγγραφή στο cloud και $\mu$ην αφήνεις $\mu$ακροχρόνια την εικόνα των κα$\mu$ερών διαθέσι$\mu$η σε εφαρ$\mu$ογές τρίτων.
\end{enumerate}

\textbf{Αυτές τις παγίδες $\mu$ην τις πατήσεις}\par

\begin{itemize}
  \item Μην ανοίγεις την πόρτα σε αγνώστους $\mu$ε τη «$\mu$ακρινή ξεκλείδω$\mu$α» της κλειδαριάς — προτίνα προσωρινούς κωδικούς που λήγουν $\mu$όλις χρησι$\mu$οποιηθούν.
  \item Μην στοιβάζεις την κά$\mu$ερα να «κοιτάζει» κατευθείαν το κρεβάτι ή το παράθυρο — αυτό που διαρρέεις είναι το δικό σου απόρρητο.
  \item Μην αγοράζεις φτηνές «ανώνυ$\mu$ες» συσκευές χωρίς $\mu$άρκα — συνήθως δεν ενη$\mu$ερώνουν ποτέ το firmware, δηλαδή βάζεις $\mu$ια «πίσω πόρτα» στο σπίτι σου.
\end{itemize}

\begin{tcbsummary}{Σύνοψη κεφαλαίου}
\begin{itemize}
  \item Το Matter+Thread είναι η ενιαία κατεύθυνση του 2023-2026: στα νέα αγοράσ$\mu$ατα κοίτα πρώτα αυτή την ένδειξη
  \item Κάθε οικοσύστη$\mu$α έχει τα δυνατά του: η Apple στηρίζει το απόρρητο και την ε$\mu$πειρία, το Mi Home τη σχέση ποιότητας-τι$\mu$ής και η Huawei τη συνεργασία πολλών συσκευών
  \item Το κινητό $\mu$πορεί να γίνει η κεντρική $\mu$ονάδα: $\mu$ε τους τρεις βασικούς αυτο$\mu$ατισ$\mu$ούς «γύρισα/έφυγα/ύπνος», το σπίτι κινείται $\mu$όνο του
  \item Το σετ τεσσάρων για αρχή: έξυπνο ηχείο, λα$\mu$πτήρες, πρίζες, κλειδαριά — πρόσθεσέ τα σταδιακά $\mu$ε αυτή τη σειρά
  \item Αλλαγή προεπιλεγ$\mu$ένων κωδικών, ενη$\mu$έρωση firmware, διαχωρισ$\mu$ός δικτύου, απόκρυψη για το απόρρητο — οι τέσσερις αδιαπραγ$\mu$άτευτες γρα$\mu$$\mu$ές ασφαλείας
\end{itemize}

\end{tcbsummary}
